\documentclass[../proyecto.tex]{subfiles}
\graphicspath{{\subfix{../images/}}}
\begin{document}

\section{Trabajo avanzado}

% máximo 1 página

\subsection{Prototipado electrónico}

Se ha realizado investigación en nuevas formas de prototipado electrónico, de automatización de procesos de fabricación de electrónica asistida por máquinas, y de documentación de procesos de diseño electrónico. Se han construido los primeros prototipos de relojes, amplificadores y secuenciadores, y se han puesto a prueba mediante su enseñanza en nuevos cursos de pregrado de diseño de interacción en Chile.

\subsection{Prototipado computacional}

Se ha realizado investigación de nuevas formas de automatización computacional, incluyendo la escritura de código y la creación automatizada de documentación, además de estrategias de publicación de bibliotecas de código. Se han ya publicado los primeros prototipos de bibliotecas y herramientas en repositorios públicos.

\subsection{Cursos de electrónica y computación}

En 2024 se crearon en conjunto con el académico y artista Matías Serrano una serie de cursos de pregrado que dictamos en la Escuela de Diseño de la Universidad Diego Portales: el primer semestre dictamos "DIS8644 Taller de diseño de máquinas electrónicas", y en el segundo semestre " DIS8645 Taller de diseño de máquinas computacionales".

La primera versión de cada curso las dictamos en el año 2025, y las segundas versiones este 2026, con la métrica de que son los talleres de tercer y cuarto año de diseño con mayor lista de espera y cantidad de estudiantes inscritos en la Escuela de Diseño de la Universidad Diego Portales.

Como parte de nuestras investigaciones doctorales y prácticas conjuntas, empezamos este 2026 a escribir un libro académico de pregrado titulado "Máquinas sonoras", que esperamos sirva de antesala a nuestras tesis doctorales individuales, y que se convierta en un texto de referencia para cursos de enseñanza de sonido, electrónica y programación como los que dictamos.

También desde 2026 se empezó a dictar \textit{DIS09214 Pensamiento computacional}, un nuevo curso obligatorio para segundo año de diseño, enfocado en los fundamentos de la computación aplicados a la creación de gráficas interactivas con la herramienta p5.js. Se están sistematizando y ampliando los apuntes del curso para su publicación como libro académico, con la esperanza de que se convierta en un aporte significativo en diseño, artes y disciplinas afines.

\subsection{Infraestructura comercial}

A nivel individual se ha realizado la infraestructura comercial y técnica, incluyendo la constitución de una empresa, Piruetas SpA, que permite financiar la investigación y desarrollo de los productos asociados a esta tesis doctoral, además de la contratación de personal, pasantes académicos y practicantes profesionales, desde el año 2022 hasta la fecha.

\end{document}